% ============================================================
% Árbol de Problemas — Causas, Problema Central y Efectos
% Proyecto BTH Innovación Tecnológica — Capítulo 2
% Formato: TikZ vectorial editable (APA 7)
% Uso: % ============================================================
% Árbol de Problemas — Causas, Problema Central y Efectos
% Proyecto BTH Innovación Tecnológica — Capítulo 2
% Formato: TikZ vectorial editable (APA 7)
% Uso: % ============================================================
% Árbol de Problemas — Causas, Problema Central y Efectos
% Proyecto BTH Innovación Tecnológica — Capítulo 2
% Formato: TikZ vectorial editable (APA 7)
% Uso: % ============================================================
% Árbol de Problemas — Causas, Problema Central y Efectos
% Proyecto BTH Innovación Tecnológica — Capítulo 2
% Formato: TikZ vectorial editable (APA 7)
% Uso: \input{imagenes/arbol_problemas} dentro de figure
% ============================================================
\begin{tikzpicture}[
    % --- Estilos base ---
    causa/.style={
        rectangle, rounded corners=3pt,
        draw=primary!75, fill=lightgray,
        text=darktext, align=center,
        font=\scriptsize,
        text width=4.2cm, minimum height=1.75cm,
        inner sep=4pt
    },
    efecto/.style={
        rectangle, rounded corners=3pt,
        draw=primary!75, fill=secondary!13,
        text=darktext, align=center,
        font=\scriptsize,
        text width=4.2cm, minimum height=1.75cm,
        inner sep=4pt
    },
    problema/.style={
        rectangle, rounded corners=5pt,
        draw=primary, fill=primary,
        text=white, align=center,
        font=\small\bfseries,
        text width=11.8cm, minimum height=2.05cm,
        inner sep=6pt
    },
    flecha/.style={
        -{Stealth[length=2.8mm,width=2.1mm]},
        thick, draw=primary!85
    },
    etiqueta/.style={
        font=\scriptsize\bfseries, text=primary, align=center
    },
    nota/.style={
        font=\tiny\itshape, text=darktext!70, align=center
    }
]

% --- EFECTOS (nivel superior) ---
\node[efecto] (E1) at (-4.9, 3.6) {Concurrencia masiva\\para reseteo manual\\y \textbf{pérdida de cupos}\\ \textbf{vocacionales}};
\node[efecto] (E2) at (0, 3.6) {Expedientes truncos,\\0\% trazabilidad y\\ \textbf{sobrecarga de}\\ \textbf{secretaría académica}};
\node[efecto] (E3) at (4.9, 3.6) {Desequilibrios por\\taller, errores de\\transcripción y\\ \textbf{cruces horarios}\\MTP $\leftrightarrow$ U.E.};

% Etiqueta EFECTOS
\node[etiqueta] (labEfectos) at (0, 5.05) {EFECTOS \quad \normalfont\tiny\textit{(manifestaciones visibles)}};
\draw[primary!40, thick] (-6.35,4.78) -- (6.35,4.78);

% --- PROBLEMA CENTRAL (nivel medio) ---
\node[problema] (P) at (0, 0.85) {Ausencia de un \textbf{sistema web integral} en el MTP San Julián BTH\\que garantice \textbf{recuperación autónoma de credenciales},\\ \textbf{matriculación documental auditada} y \textbf{asignación automatizada}\\ a cursos y paralelos por \textbf{unidad educativa de procedencia}};

% Etiqueta PROBLEMA CENTRAL lateral
\node[etiqueta, rotate=90, anchor=center, font=\scriptsize\bfseries, text=primary!90] at (-7.15, 0.85) {PROBLEMA CENTRAL};

% --- CAUSAS RAÍCES (nivel inferior) ---
\node[causa] (C1) at (-4.9, -1.95) {Plataforma V1.0 (2025)\\ \textbf{sin autoservicio}\\de credenciales\\(sin validación \textbf{C.I. +}\\ \textbf{fecha nacimiento})};
\node[causa] (C2) at (0, -1.95) {Inexistencia de \textbf{módulo}\\ \textbf{de matriculación}\\ \textbf{documental} en BD\\(0\% trazabilidad)};
\node[causa] (C3) at (4.9, -1.95) {Asignación \textbf{manual}\\en planillas externas\\por \textbf{colegio de origen}\\(U.E. de procedencia)};

% Etiqueta CAUSAS
\node[etiqueta] (labCausas) at (0, -3.55) {CAUSAS RAÍCES \quad \normalfont\tiny\textit{(origen técnico-funcional)}};
\draw[primary!40, thick] (-6.35,-3.28) -- (6.35,-3.28);

% --- FLECHAS: Causas → Problema Central ---
\draw[flecha] (C1.north) -- (-2.8, -0.20);
\draw[flecha] (C2.north) -- (0, -0.20);
\draw[flecha] (C3.north) -- (2.8, -0.20);

% --- FLECHAS: Problema Central → Efectos ---
\draw[flecha] (-2.8, 1.92) -- (E1.south);
\draw[flecha] (0, 1.92) -- (E2.south);
\draw[flecha] (2.8, 1.92) -- (E3.south);

% --- Contexto inferior (población/diagnóstico) ---
\node[nota] at (0, -4.05) {Contexto: $\sim$400 postulantes/año $\cdot$ 7 especialidades $\cdot$ Diagnóstico 2022--2025 + relevamiento 1er bimestre 2026 (60 estudiantes + 5 funcionarios)};

% --- Marco sutil opcional (desactivado, respeta APA) ---
% \begin{scope}[on background layer]
%   \node[fit=(E1)(E3)(labEfectos), draw=primary!18, rounded corners=6pt, inner sep=6pt] {};
%   \node[fit=(C1)(C3)(labCausas), draw=primary!18, rounded corners=6pt, inner sep=6pt] {};
% \end{scope}

\end{tikzpicture}
 dentro de figure
% ============================================================
\begin{tikzpicture}[
    % --- Estilos base ---
    causa/.style={
        rectangle, rounded corners=3pt,
        draw=primary!75, fill=lightgray,
        text=darktext, align=center,
        font=\scriptsize,
        text width=4.2cm, minimum height=1.75cm,
        inner sep=4pt
    },
    efecto/.style={
        rectangle, rounded corners=3pt,
        draw=primary!75, fill=secondary!13,
        text=darktext, align=center,
        font=\scriptsize,
        text width=4.2cm, minimum height=1.75cm,
        inner sep=4pt
    },
    problema/.style={
        rectangle, rounded corners=5pt,
        draw=primary, fill=primary,
        text=white, align=center,
        font=\small\bfseries,
        text width=11.8cm, minimum height=2.05cm,
        inner sep=6pt
    },
    flecha/.style={
        -{Stealth[length=2.8mm,width=2.1mm]},
        thick, draw=primary!85
    },
    etiqueta/.style={
        font=\scriptsize\bfseries, text=primary, align=center
    },
    nota/.style={
        font=\tiny\itshape, text=darktext!70, align=center
    }
]

% --- EFECTOS (nivel superior) ---
\node[efecto] (E1) at (-4.9, 3.6) {Concurrencia masiva\\para reseteo manual\\y \textbf{pérdida de cupos}\\ \textbf{vocacionales}};
\node[efecto] (E2) at (0, 3.6) {Expedientes truncos,\\0\% trazabilidad y\\ \textbf{sobrecarga de}\\ \textbf{secretaría académica}};
\node[efecto] (E3) at (4.9, 3.6) {Desequilibrios por\\taller, errores de\\transcripción y\\ \textbf{cruces horarios}\\MTP $\leftrightarrow$ U.E.};

% Etiqueta EFECTOS
\node[etiqueta] (labEfectos) at (0, 5.05) {EFECTOS \quad \normalfont\tiny\textit{(manifestaciones visibles)}};
\draw[primary!40, thick] (-6.35,4.78) -- (6.35,4.78);

% --- PROBLEMA CENTRAL (nivel medio) ---
\node[problema] (P) at (0, 0.85) {Ausencia de un \textbf{sistema web integral} en el MTP San Julián BTH\\que garantice \textbf{recuperación autónoma de credenciales},\\ \textbf{matriculación documental auditada} y \textbf{asignación automatizada}\\ a cursos y paralelos por \textbf{unidad educativa de procedencia}};

% Etiqueta PROBLEMA CENTRAL lateral
\node[etiqueta, rotate=90, anchor=center, font=\scriptsize\bfseries, text=primary!90] at (-7.15, 0.85) {PROBLEMA CENTRAL};

% --- CAUSAS RAÍCES (nivel inferior) ---
\node[causa] (C1) at (-4.9, -1.95) {Plataforma V1.0 (2025)\\ \textbf{sin autoservicio}\\de credenciales\\(sin validación \textbf{C.I. +}\\ \textbf{fecha nacimiento})};
\node[causa] (C2) at (0, -1.95) {Inexistencia de \textbf{módulo}\\ \textbf{de matriculación}\\ \textbf{documental} en BD\\(0\% trazabilidad)};
\node[causa] (C3) at (4.9, -1.95) {Asignación \textbf{manual}\\en planillas externas\\por \textbf{colegio de origen}\\(U.E. de procedencia)};

% Etiqueta CAUSAS
\node[etiqueta] (labCausas) at (0, -3.55) {CAUSAS RAÍCES \quad \normalfont\tiny\textit{(origen técnico-funcional)}};
\draw[primary!40, thick] (-6.35,-3.28) -- (6.35,-3.28);

% --- FLECHAS: Causas → Problema Central ---
\draw[flecha] (C1.north) -- (-2.8, -0.20);
\draw[flecha] (C2.north) -- (0, -0.20);
\draw[flecha] (C3.north) -- (2.8, -0.20);

% --- FLECHAS: Problema Central → Efectos ---
\draw[flecha] (-2.8, 1.92) -- (E1.south);
\draw[flecha] (0, 1.92) -- (E2.south);
\draw[flecha] (2.8, 1.92) -- (E3.south);

% --- Contexto inferior (población/diagnóstico) ---
\node[nota] at (0, -4.05) {Contexto: $\sim$400 postulantes/año $\cdot$ 7 especialidades $\cdot$ Diagnóstico 2022--2025 + relevamiento 1er bimestre 2026 (60 estudiantes + 5 funcionarios)};

% --- Marco sutil opcional (desactivado, respeta APA) ---
% \begin{scope}[on background layer]
%   \node[fit=(E1)(E3)(labEfectos), draw=primary!18, rounded corners=6pt, inner sep=6pt] {};
%   \node[fit=(C1)(C3)(labCausas), draw=primary!18, rounded corners=6pt, inner sep=6pt] {};
% \end{scope}

\end{tikzpicture}
 dentro de figure
% ============================================================
\begin{tikzpicture}[
    % --- Estilos base ---
    causa/.style={
        rectangle, rounded corners=3pt,
        draw=primary!75, fill=lightgray,
        text=darktext, align=center,
        font=\scriptsize,
        text width=4.2cm, minimum height=1.75cm,
        inner sep=4pt
    },
    efecto/.style={
        rectangle, rounded corners=3pt,
        draw=primary!75, fill=secondary!13,
        text=darktext, align=center,
        font=\scriptsize,
        text width=4.2cm, minimum height=1.75cm,
        inner sep=4pt
    },
    problema/.style={
        rectangle, rounded corners=5pt,
        draw=primary, fill=primary,
        text=white, align=center,
        font=\small\bfseries,
        text width=11.8cm, minimum height=2.05cm,
        inner sep=6pt
    },
    flecha/.style={
        -{Stealth[length=2.8mm,width=2.1mm]},
        thick, draw=primary!85
    },
    etiqueta/.style={
        font=\scriptsize\bfseries, text=primary, align=center
    },
    nota/.style={
        font=\tiny\itshape, text=darktext!70, align=center
    }
]

% --- EFECTOS (nivel superior) ---
\node[efecto] (E1) at (-4.9, 3.6) {Concurrencia masiva\\para reseteo manual\\y \textbf{pérdida de cupos}\\ \textbf{vocacionales}};
\node[efecto] (E2) at (0, 3.6) {Expedientes truncos,\\0\% trazabilidad y\\ \textbf{sobrecarga de}\\ \textbf{secretaría académica}};
\node[efecto] (E3) at (4.9, 3.6) {Desequilibrios por\\taller, errores de\\transcripción y\\ \textbf{cruces horarios}\\MTP $\leftrightarrow$ U.E.};

% Etiqueta EFECTOS
\node[etiqueta] (labEfectos) at (0, 5.05) {EFECTOS \quad \normalfont\tiny\textit{(manifestaciones visibles)}};
\draw[primary!40, thick] (-6.35,4.78) -- (6.35,4.78);

% --- PROBLEMA CENTRAL (nivel medio) ---
\node[problema] (P) at (0, 0.85) {Ausencia de un \textbf{sistema web integral} en el MTP San Julián BTH\\que garantice \textbf{recuperación autónoma de credenciales},\\ \textbf{matriculación documental auditada} y \textbf{asignación automatizada}\\ a cursos y paralelos por \textbf{unidad educativa de procedencia}};

% Etiqueta PROBLEMA CENTRAL lateral
\node[etiqueta, rotate=90, anchor=center, font=\scriptsize\bfseries, text=primary!90] at (-7.15, 0.85) {PROBLEMA CENTRAL};

% --- CAUSAS RAÍCES (nivel inferior) ---
\node[causa] (C1) at (-4.9, -1.95) {Plataforma V1.0 (2025)\\ \textbf{sin autoservicio}\\de credenciales\\(sin validación \textbf{C.I. +}\\ \textbf{fecha nacimiento})};
\node[causa] (C2) at (0, -1.95) {Inexistencia de \textbf{módulo}\\ \textbf{de matriculación}\\ \textbf{documental} en BD\\(0\% trazabilidad)};
\node[causa] (C3) at (4.9, -1.95) {Asignación \textbf{manual}\\en planillas externas\\por \textbf{colegio de origen}\\(U.E. de procedencia)};

% Etiqueta CAUSAS
\node[etiqueta] (labCausas) at (0, -3.55) {CAUSAS RAÍCES \quad \normalfont\tiny\textit{(origen técnico-funcional)}};
\draw[primary!40, thick] (-6.35,-3.28) -- (6.35,-3.28);

% --- FLECHAS: Causas → Problema Central ---
\draw[flecha] (C1.north) -- (-2.8, -0.20);
\draw[flecha] (C2.north) -- (0, -0.20);
\draw[flecha] (C3.north) -- (2.8, -0.20);

% --- FLECHAS: Problema Central → Efectos ---
\draw[flecha] (-2.8, 1.92) -- (E1.south);
\draw[flecha] (0, 1.92) -- (E2.south);
\draw[flecha] (2.8, 1.92) -- (E3.south);

% --- Contexto inferior (población/diagnóstico) ---
\node[nota] at (0, -4.05) {Contexto: $\sim$400 postulantes/año $\cdot$ 7 especialidades $\cdot$ Diagnóstico 2022--2025 + relevamiento 1er bimestre 2026 (60 estudiantes + 5 funcionarios)};

% --- Marco sutil opcional (desactivado, respeta APA) ---
% \begin{scope}[on background layer]
%   \node[fit=(E1)(E3)(labEfectos), draw=primary!18, rounded corners=6pt, inner sep=6pt] {};
%   \node[fit=(C1)(C3)(labCausas), draw=primary!18, rounded corners=6pt, inner sep=6pt] {};
% \end{scope}

\end{tikzpicture}
 dentro de figure
% ============================================================
\begin{tikzpicture}[
    % --- Estilos base ---
    causa/.style={
        rectangle, rounded corners=3pt,
        draw=primary!75, fill=lightgray,
        text=darktext, align=center,
        font=\scriptsize,
        text width=4.2cm, minimum height=1.75cm,
        inner sep=4pt
    },
    efecto/.style={
        rectangle, rounded corners=3pt,
        draw=primary!75, fill=secondary!13,
        text=darktext, align=center,
        font=\scriptsize,
        text width=4.2cm, minimum height=1.75cm,
        inner sep=4pt
    },
    problema/.style={
        rectangle, rounded corners=5pt,
        draw=primary, fill=primary,
        text=white, align=center,
        font=\small\bfseries,
        text width=11.8cm, minimum height=2.05cm,
        inner sep=6pt
    },
    flecha/.style={
        -{Stealth[length=2.8mm,width=2.1mm]},
        thick, draw=primary!85
    },
    etiqueta/.style={
        font=\scriptsize\bfseries, text=primary, align=center
    },
    nota/.style={
        font=\tiny\itshape, text=darktext!70, align=center
    }
]

% --- EFECTOS (nivel superior) ---
\node[efecto] (E1) at (-4.9, 3.6) {Concurrencia masiva\\para reseteo manual\\y \textbf{pérdida de cupos}\\ \textbf{vocacionales}};
\node[efecto] (E2) at (0, 3.6) {Expedientes truncos,\\0\% trazabilidad y\\ \textbf{sobrecarga de}\\ \textbf{secretaría académica}};
\node[efecto] (E3) at (4.9, 3.6) {Desequilibrios por\\taller, errores de\\transcripción y\\ \textbf{cruces horarios}\\MTP $\leftrightarrow$ U.E.};

% Etiqueta EFECTOS
\node[etiqueta] (labEfectos) at (0, 5.05) {EFECTOS \quad \normalfont\tiny\textit{(manifestaciones visibles)}};
\draw[primary!40, thick] (-6.35,4.78) -- (6.35,4.78);

% --- PROBLEMA CENTRAL (nivel medio) ---
\node[problema] (P) at (0, 0.85) {Ausencia de un \textbf{sistema web integral} en el MTP San Julián BTH\\que garantice \textbf{recuperación autónoma de credenciales},\\ \textbf{matriculación documental auditada} y \textbf{asignación automatizada}\\ a cursos y paralelos por \textbf{unidad educativa de procedencia}};

% Etiqueta PROBLEMA CENTRAL lateral
\node[etiqueta, rotate=90, anchor=center, font=\scriptsize\bfseries, text=primary!90] at (-7.15, 0.85) {PROBLEMA CENTRAL};

% --- CAUSAS RAÍCES (nivel inferior) ---
\node[causa] (C1) at (-4.9, -1.95) {Plataforma V1.0 (2025)\\ \textbf{sin autoservicio}\\de credenciales\\(sin validación \textbf{C.I. +}\\ \textbf{fecha nacimiento})};
\node[causa] (C2) at (0, -1.95) {Inexistencia de \textbf{módulo}\\ \textbf{de matriculación}\\ \textbf{documental} en BD\\(0\% trazabilidad)};
\node[causa] (C3) at (4.9, -1.95) {Asignación \textbf{manual}\\en planillas externas\\por \textbf{colegio de origen}\\(U.E. de procedencia)};

% Etiqueta CAUSAS
\node[etiqueta] (labCausas) at (0, -3.55) {CAUSAS RAÍCES \quad \normalfont\tiny\textit{(origen técnico-funcional)}};
\draw[primary!40, thick] (-6.35,-3.28) -- (6.35,-3.28);

% --- FLECHAS: Causas → Problema Central ---
\draw[flecha] (C1.north) -- (-2.8, -0.20);
\draw[flecha] (C2.north) -- (0, -0.20);
\draw[flecha] (C3.north) -- (2.8, -0.20);

% --- FLECHAS: Problema Central → Efectos ---
\draw[flecha] (-2.8, 1.92) -- (E1.south);
\draw[flecha] (0, 1.92) -- (E2.south);
\draw[flecha] (2.8, 1.92) -- (E3.south);

% --- Contexto inferior (población/diagnóstico) ---
\node[nota] at (0, -4.05) {Contexto: $\sim$400 postulantes/año $\cdot$ 7 especialidades $\cdot$ Diagnóstico 2022--2025 + relevamiento 1er bimestre 2026 (60 estudiantes + 5 funcionarios)};

% --- Marco sutil opcional (desactivado, respeta APA) ---
% \begin{scope}[on background layer]
%   \node[fit=(E1)(E3)(labEfectos), draw=primary!18, rounded corners=6pt, inner sep=6pt] {};
%   \node[fit=(C1)(C3)(labCausas), draw=primary!18, rounded corners=6pt, inner sep=6pt] {};
% \end{scope}

\end{tikzpicture}
